\documentclass[12pt,a4paper]{article}

% ===== PACKAGES =====
\usepackage[utf8]{inputenc}
\usepackage[T1]{fontenc}
\usepackage{times}
\usepackage[margin=0.9in, top=1in, bottom=0.9in]{geometry}
\usepackage{graphicx}
\usepackage{booktabs}
\usepackage{array}
\usepackage{longtable}
\usepackage{hyperref}
\usepackage{xcolor}
\usepackage{titlesec}
\usepackage{fancyhdr}
\usepackage{setspace}
\usepackage{parskip}
\usepackage{enumitem}
\usepackage{amsmath}
\usepackage{float}
\usepackage{textcomp}

% ===== COLORS =====
\definecolor{navyblue}{HTML}{003366}
\definecolor{lightgray}{HTML}{F5F5F5}

% ===== HYPERREF SETUP =====
\hypersetup{
    colorlinks=true,
    linkcolor=navyblue,
    urlcolor=navyblue,
    citecolor=navyblue
}

% ===== SECTION FORMATTING (Compact) =====
\titleformat{\section}
    {\normalfont\large\bfseries\color{navyblue}}
    {\thesection.}{0.5em}{}
\titleformat{\subsection}
    {\normalfont\normalsize\bfseries\color{navyblue}}
    {\thesubsection}{0.5em}{}
\titleformat{\subsubsection}
    {\normalfont\small\bfseries}
    {\thesubsubsection}{0.5em}{}

% Compact section spacing
\titlespacing*{\section}{0pt}{10pt plus 2pt minus 2pt}{6pt plus 1pt minus 1pt}
\titlespacing*{\subsection}{0pt}{8pt plus 2pt minus 2pt}{4pt plus 1pt minus 1pt}

% ===== HEADER/FOOTER =====
\pagestyle{fancy}
\fancyhf{}
\setlength{\headheight}{14pt}
\fancyhead[C]{\small MTech Progress Report - Semester III}
\fancyfoot[C]{\thepage}
\renewcommand{\headrulewidth}{0.4pt}

% ===== LINE SPACING (Single for compactness) =====
\setstretch{1.15}

% ===== DOCUMENT BEGIN =====
\begin{document}

% ========================================
% PAGE 1: COVER PAGE
% ========================================
\thispagestyle{empty}
\begin{center}

\vspace*{0.5cm}

% IIT Patna Logo
\includegraphics[width=3.5cm]{iitp-logo.png}

\vspace{0.5cm}

% IIT Patna Header
{\Large\bfseries INDIAN INSTITUTE OF TECHNOLOGY PATNA}\\[0.2cm]
{\normalsize Department of Artificial Intelligence and Data Science Engineering}

\vspace{0.8cm}

% Decorative line
\noindent\rule{0.8\textwidth}{1pt}

\vspace{0.6cm}

{\Large\bfseries SEMESTER III PROGRESS REPORT}

\vspace{0.6cm}

\noindent\rule{0.8\textwidth}{1pt}

\vspace{1cm}

% Title Box
\fbox{\parbox{0.85\textwidth}{\centering
\vspace{0.4cm}
{\large\bfseries Real-Time Bidirectional Continuous Sign Language\\[0.2cm] Interpretation Using Deep Learning}
\vspace{0.4cm}
}}

\vspace{1.2cm}

% Student Details
{\normalsize\bfseries Submitted by:}\\[0.4cm]
{\Large\bfseries ARCHIE NARAYAN}\\[0.2cm]
{\normalsize Roll No: 24A03RES119}\\[0.2cm]
{\normalsize MTech (Artificial Intelligence and Data Science Engineering)}

\vspace{1.2cm}

\noindent\rule{0.8\textwidth}{1pt}

\vspace{0.6cm}

{\normalsize\bfseries December 2025}

\vfill

{\normalsize Indian Institute of Technology Patna\\
Bihta, Patna - 801106, Bihar, India}

\end{center}

\newpage

% ========================================
% PAGE 2: ABSTRACT & INTRODUCTION
% ========================================

\section{Abstract}

Sign language serves as the primary mode of communication for over 70 million deaf and hard-of-hearing individuals worldwide. However, the communication barrier between sign language users and the hearing population remains a significant challenge. This project presents a comprehensive deep learning-based system for \textbf{bidirectional} continuous sign language interpretation, enabling seamless communication in both directions.

The system comprises two main pipelines: (1) \textbf{Sign Language Recognition (SLR)} -- converting continuous sign language videos to text, and (2) \textbf{Sign Language Production (SLP)} -- converting text/speech to sign language videos. For the SLR component, we implemented a \textbf{Hybrid CTC+Attention Transformer architecture} trained on the RWTH-PHOENIX-Weather 2014 dataset, achieving a Word Error Rate (WER) of \textbf{50.91\%} on the test set. The architecture combines a pretrained ResNet-18 CNN backbone with a 6-layer Transformer encoder, using joint CTC and cross-entropy losses to prevent the common CTC collapse problem.

For the SLP component, we developed a retrieval-based video synthesis system that translates text to gloss sequences and retrieves corresponding sign video segments from an indexed database, enabling real-time text-to-sign conversion.

\textbf{Keywords:} Sign Language Recognition, Continuous Sign Language, CTC, Transformer, Attention Mechanism, Deep Learning, PHOENIX-2014

\section{Introduction}

\subsection{Motivation}

According to the World Health Organization, approximately 466 million people worldwide have disabling hearing loss. Sign language is a complete, natural language with its own grammar and syntax, used by deaf communities across the globe. However, the lack of sign language interpreters creates significant barriers in education, healthcare, employment, and social interactions.

Advances in deep learning and computer vision have opened new possibilities for automated sign language interpretation. Unlike isolated sign recognition (recognizing individual signs), \textbf{Continuous Sign Language Recognition (CSLR)} deals with connected signs in natural sentences, presenting unique challenges including:

\begin{itemize}[noitemsep]
    \item Variable-length sequences with no explicit word boundaries
    \item Co-articulation effects between consecutive signs
    \item Signer variations in speed, style, and appearance
    \item Limited availability of annotated datasets
\end{itemize}

\subsection{Objectives}

The primary objectives of this project are:

\begin{enumerate}[noitemsep]
    \item \textbf{Develop a robust CSLR system} capable of recognizing continuous sign language from video input with competitive accuracy on benchmark datasets.
    
    \item \textbf{Implement a bidirectional interpretation system} that enables communication in both directions -- from sign language to text and from text/speech to sign language.
    
    \item \textbf{Create real-time demonstration systems} including camera-based sign recognition and text-to-sign video generation.
    
    \item \textbf{Document technical learnings} regarding deep learning architectures for sequence-to-sequence tasks with variable-length alignment.
\end{enumerate}

\section{Literature Review}

\subsection{Evolution of Sign Language Recognition}

Sign language recognition has evolved significantly over the past decade. Early approaches relied on hand-crafted features and Hidden Markov Models (HMMs). The introduction of deep learning, particularly Convolutional Neural Networks (CNNs) and Recurrent Neural Networks (RNNs), marked a paradigm shift in the field.

\begin{table}[H]
\centering
\caption{Comparison of State-of-the-Art Methods on PHOENIX-2014}
\label{tab:sota}
\begin{tabular}{@{}lllcc@{}}
\toprule
\textbf{Method} & \textbf{Year} & \textbf{Architecture} & \textbf{Dev WER} & \textbf{Test WER} \\
\midrule
Koller et al. & 2017 & CNN-HMM-RNN & 26.0\% & 26.8\% \\
Camgöz et al. & 2020 & Transformer + CTC & 21.8\% & 24.5\% \\
Niu \& Mak & 2020 & Stochastic CSLR & 25.4\% & 24.0\% \\
Zhu et al. (VAC) & 2021 & Visual Alignment & 19.3\% & 21.2\% \\
CorrNet & 2023 & Correlation Net & 17.3\% & 17.8\% \\
\textbf{Ours} & \textbf{2025} & \textbf{Hybrid CTC+Attention} & \textbf{51.44\%} & \textbf{50.91\%} \\
\bottomrule
\end{tabular}
\end{table}

\subsection{CTC-Based Approaches}

Connectionist Temporal Classification (CTC) introduced by Graves et al. (2006) enables training sequence-to-sequence models without requiring frame-level alignments. CTC has been widely adopted in speech recognition and sign language recognition due to its ability to handle variable-length sequences.

\subsection{Hybrid CTC+Attention Models}

Recent works have shown that pure CTC models are prone to collapse, outputting only blank tokens. The hybrid CTC+Attention approach (Watanabe et al., 2017) combines the alignment capabilities of CTC with the sequence modeling power of attention-based decoders, providing more stable training.

\section{Problem Statement}

\subsection{Technical Challenges}

The development of a practical CSLR system faces several technical challenges:

\begin{enumerate}[noitemsep]
    \item \textbf{CTC Collapse Problem:} Pure CTC models tend to learn trivial solutions by outputting only blank tokens, resulting in 100\% word error rate regardless of input.
    
    \item \textbf{Long-Range Dependencies:} Sign language sentences can span hundreds of frames, requiring models to capture long-range temporal dependencies.
    
    \item \textbf{Limited Labeled Data:} Unlike speech recognition, sign language datasets are relatively small, making deep learning approaches prone to overfitting.
    
    \item \textbf{Variable Speed and Duration:} Signers perform signs at different speeds, and sign duration varies based on context and signer preference.
\end{enumerate}

\subsection{Research Questions}

This project addresses the following research questions:

\textbf{RQ1:} How can we prevent CTC collapse in continuous sign language recognition while maintaining the benefits of CTC-based alignment?

\textbf{RQ2:} What is the optimal architecture for capturing both short-term motion patterns and long-term semantic dependencies in sign language videos?

\textbf{RQ3:} How can we create a practical bidirectional system that enables real-time interpretation?

\section{Methodology}

\subsection{Dataset}

We use the \textbf{RWTH-PHOENIX-Weather 2014} dataset, the standard benchmark for continuous sign language recognition:

\begin{table}[H]
\centering
\caption{RWTH-PHOENIX-Weather 2014 Dataset Statistics}
\label{tab:dataset}
\begin{tabular}{@{}ll@{}}
\toprule
\textbf{Property} & \textbf{Value} \\
\midrule
Language & German Sign Language (DGS) \\
Domain & Weather forecasts \\
Vocabulary & 1,236 unique glosses \\
Training samples & 5,672 \\
Development samples & 540 \\
Test samples & 629 \\
Frame dimensions & 210 $\times$ 260 pixels \\
Frame rate & 25 FPS \\
Total size & $\sim$53 GB \\
\bottomrule
\end{tabular}
\end{table}

\subsection{System Architecture}

Our system employs a \textbf{Hybrid CTC+Attention architecture} consisting of four main components:

\textbf{Architecture Flow:} Video Frames $\rightarrow$ \textbf{ResNet-18 CNN} (512-dim features per frame) $\rightarrow$ \textbf{Transformer Encoder} (6 layers, 8 heads, d\_model=512) $\rightarrow$ Dual heads: \textbf{CTC Head} (0.3 weight) + \textbf{Attention Decoder} (0.7 weight, 3 layers) $\rightarrow$ Joint Loss.

\subsection{Loss Functions}

\textbf{CTC Loss:} Enables alignment-free training by marginalizing over all possible alignments between input frames and output glosses.

\textbf{Cross-Entropy Loss:} Applied to the decoder output with label smoothing (0.1) to improve generalization.

\textbf{Joint Loss:} The weighted combination prevents CTC collapse while maintaining alignment benefits:
\begin{equation}
L_{total} = \lambda_{CTC} \times L_{CTC} + \lambda_{CE} \times L_{CE}
\end{equation}
where $\lambda_{CTC} = 0.3$ and $\lambda_{CE} = 0.7$.

\subsection{Training Strategy}

\begin{itemize}[noitemsep]
    \item \textbf{Optimizer:} AdamW with weight decay 0.05
    \item \textbf{Learning Rate:} $1 \times 10^{-4}$ with warmup (5 epochs) + cosine annealing
    \item \textbf{Batch Size:} 4 (limited by GPU memory)
    \item \textbf{Max Frames:} 64 frames per video
    \item \textbf{Gradient Clipping:} max\_norm = 5.0
    \item \textbf{Early Stopping:} Patience of 15 epochs
\end{itemize}

\section{Implementation Details}

\subsection{Environment \& Configuration}

\textbf{Environment:} Python 3.10, PyTorch 2.0, NVIDIA GPU with CUDA, $\sim$20 hours training (100 epochs).

\textbf{Model:} ResNet-18 backbone (ImageNet pretrained), d\_model=512, 6 encoder layers, 3 decoder layers, 8 attention heads, feed-forward dim=2048, dropout=0.2, total 45.2M parameters.

\textbf{Augmentation:} Spatial (random crop 0.85-1.0, rotation $\pm 8^{\circ}$, brightness/contrast jitter), Temporal (speed 0.8-1.2x, uniform sampling). \textit{Note: No horizontal flip} (changes sign meaning).

\section{Experimental Results}

\subsection{Main Results}

\begin{table}[H]
\centering
\caption{Final Model Performance}
\begin{tabular}{@{}lcc@{}}
\toprule
\textbf{Metric} & \textbf{Development Set} & \textbf{Test Set} \\
\midrule
Word Error Rate (WER) & 51.44\% & \textbf{50.91\%} \\
Best Epoch & 98 & -- \\
Training Loss & 0.0006 & -- \\
\bottomrule
\end{tabular}
\end{table}

\subsection{Training Progress \& CTC Collapse Prevention}

\textbf{WER Progression:} Epoch 1: 91.75\% $\rightarrow$ Epoch 10: 70.25\% $\rightarrow$ Epoch 20: 57.55\% $\rightarrow$ Epoch 40: 55.56\% $\rightarrow$ Epoch 60: 55.04\% $\rightarrow$ Epoch 80: 53.76\% $\rightarrow$ \textbf{Epoch 98: 52.85\% (Best)}.

\textbf{CTC Collapse Prevention:} Pure CTC collapses to 100\% WER (outputting only blank tokens). Our Hybrid CTC+Attention approach achieved 50.91\% WER, successfully preventing collapse through the attention decoder's stable gradients.

\subsection{Key Findings}

\begin{enumerate}[noitemsep]
    \item \textbf{Hybrid architecture prevents collapse:} The attention decoder provides stable gradients even when CTC starts to collapse.
    
    \item \textbf{Lower CTC weight is crucial:} CTC weight $> 0.5$ leads to collapse; 0.3 provides optimal balance.
    
    \item \textbf{Pretrained backbone essential:} ImageNet-pretrained ResNet-18 significantly outperforms random initialization.
    
    \item \textbf{Label smoothing helps:} 0.1 label smoothing improves generalization by $\sim$2\% WER.
\end{enumerate}

\section{Comprehensive System Implementation}

Beyond the core recognition model, we developed a \textbf{complete bidirectional system} with production-ready components:

\subsection{Data Augmentation Pipeline (Implemented)}

Custom video augmentation module preserving sign semantics:
\begin{itemize}[noitemsep]
    \item \textbf{Spatial:} Random crop (0.8-1.0 scale), rotation ($\pm 10^{\circ}$), color jitter (brightness, contrast, saturation)
    \item \textbf{Temporal:} Frame dropout, temporal masking, speed perturbation (0.8-1.2x)
    \item \textbf{Critical constraint:} No horizontal flip (inverts sign meaning)
\end{itemize}

\subsection{Speech-to-Sign Pipeline (Implemented)}

Complete reverse pipeline (\texttt{SpeechToSignPipeline} class):
\begin{enumerate}[noitemsep]
    \item \textbf{ASR Module:} Whisper integration for speech-to-text
    \item \textbf{Text-to-Gloss:} Rule-based + neural translation (Transformer-based)
    \item \textbf{Gloss-to-Video:} Retrieval-based synthesis with indexed database of 1,236 glosses
\end{enumerate}

\subsection{Real-Time Demo Systems (Implemented)}

\begin{itemize}[noitemsep]
    \item \textbf{Camera Demo:} Live sign recognition with English translation overlay
    \item \textbf{Video Demo:} Process pre-recorded sign language videos
    \item \textbf{Text-to-Sign:} Generate sign videos from text input
    \item \textbf{Streamlit Dashboard:} Training monitoring with live metrics visualization
\end{itemize}

\subsection{I3D Feature Extraction (Implemented)}

Pre-extracted I3D features from Kinetics-pretrained model for faster training:
\begin{itemize}[noitemsep]
    \item Reduces training time from 20 hours to $\sim$6 hours
    \item Enables batch size 8 (vs. 4 for end-to-end)
    \item Captures spatiotemporal patterns at multiple scales
\end{itemize}

\section{Future Work and Optimizations}

\subsection{Performance Optimizations}

Several promising directions exist to improve recognition accuracy:

\begin{enumerate}[noitemsep]
    \item \textbf{Visual Alignment Constraint (VAC):} Implement auxiliary loss enforcing monotonic alignment between video frames and gloss sequences. This addresses temporal misalignment issues and is expected to reduce WER by 3-5\%.
    
    \item \textbf{Extended Temporal Context:} Increase max\_frames from 64 to 150+ using gradient checkpointing to handle longer sign sequences. Expected improvement: 5-8\% WER reduction.
    
    \item \textbf{Self-Distillation:} Employ teacher-student framework where the model's own predictions serve as soft targets, providing additional supervision signal. Expected improvement: 2-3\% WER reduction.
    
    \item \textbf{Multi-Scale Temporal Modeling:} Add temporal convolutions at multiple scales (3, 5, 7, 11 frames) to capture both short-term motion patterns and long-term dependencies simultaneously.
    
    \item \textbf{Advanced Data Augmentation:} Implement CutMix for videos, temporal segment swapping, and adversarial training to improve robustness to signer variations.
\end{enumerate}

\subsection{Cross-Language Generalization}

\subsubsection{American Sign Language (ASL) Application}

A critical next step is validating the architecture on American Sign Language datasets:

\begin{itemize}[noitemsep]
    \item \textbf{WLASL Dataset:} Large-scale ASL dataset with 2,000+ signs and 21,000+ videos. Testing on WLASL would demonstrate cross-language transferability.
    
    \item \textbf{How2Sign Dataset:} English-ASL parallel corpus with instructional videos, enabling evaluation of the bidirectional pipeline on a different sign language.
    
    \item \textbf{Transfer Learning Strategy:} Fine-tune the PHOENIX-pretrained model on ASL data, leveraging learned spatiotemporal representations while adapting to ASL-specific patterns.
    
    \item \textbf{Cross-Language Analysis:} Compare recognition performance between DGS (German Sign Language) and ASL to identify language-specific challenges and universal features.
\end{itemize}

\subsubsection{Other Sign Languages}

\begin{itemize}[noitemsep]
    \item \textbf{Chinese Sign Language (CSL):} Evaluate on CSL-Daily dataset to test generalization to Asian sign languages with different linguistic structures.
    
    \item \textbf{British Sign Language (BSL):} Test on BSL datasets to assess performance across European sign language variants.
\end{itemize}

\subsection{Model Optimization and Deployment}

\begin{enumerate}[noitemsep]
    \item \textbf{Model Compression:} Quantization (INT8), pruning, and knowledge distillation to reduce model size for mobile deployment.
    
    \item \textbf{Inference Acceleration:} ONNX/TensorRT export, optimized attention mechanisms, and frame-level caching for real-time performance.
    
    \item \textbf{Production Deployment:} Web application with Flask/FastAPI backend and React frontend for public accessibility.
    
    \item \textbf{Mobile Application:} Native iOS/Android apps with on-device inference for offline sign language interpretation.
\end{enumerate}

\section{Conclusion}

This project established a \textbf{comprehensive foundation} for bidirectional continuous sign language interpretation using deep learning. Key contributions include:

\subsection{Technical Contributions}

\begin{enumerate}[noitemsep]
    \item \textbf{Novel Hybrid Architecture:} Solved the critical CTC collapse problem through joint CTC+Attention training, documented with comprehensive ablation studies demonstrating the necessity of hybrid approach.
    
    \item \textbf{End-to-End Bidirectional System:} Complete pipeline for sign$\leftrightarrow$text translation, integrating recognition and production in a unified framework.
    
    \item \textbf{Production-Ready Components:} Comprehensive data augmentation pipeline, I3D feature extraction, real-time demo systems, and training monitoring dashboard.
    
    \item \textbf{Reproducible Research:} Well-documented codebase with modular architecture, detailed training configurations, and extensive technical documentation.
\end{enumerate}

\subsection{Research Impact}

The work demonstrates significant progress toward practical sign language interpretation systems with clear technical novelty:
\begin{itemize}[noitemsep]
    \item Addresses the fundamental CTC collapse problem in continuous sign language recognition
    \item Provides the first complete bidirectional pipeline combining recognition and production
    \item Establishes baseline results (50.91\% WER) on PHOENIX-2014 benchmark
    \item Creates foundation for cross-language generalization to ASL and other sign languages
\end{itemize}

Future work will focus on performance optimization through VAC and self-distillation, cross-language validation on ASL datasets, and deployment optimization for real-world applications.

\section*{References}
\addcontentsline{toc}{section}{References}

\begin{enumerate}[label={[\arabic*]}, leftmargin=*, noitemsep]
    \item O. Koller, J. Forster, and H. Ney, ``Continuous sign language recognition: Towards large vocabulary statistical recognition systems handling multiple signers,'' \textit{Computer Vision and Image Understanding}, vol. 141, pp. 108-125, 2015.
    
    \item N. C. Camgöz, S. Hadfield, O. Koller, H. Ney, and R. Bowden, ``Sign Language Transformers: Joint End-to-end Sign Language Recognition and Translation,'' in \textit{Proceedings of the IEEE/CVF Conference on Computer Vision and Pattern Recognition (CVPR)}, 2020, pp. 10023-10033.
    
    \item A. Graves, S. Fernández, F. Gomez, and J. Schmidhuber, ``Connectionist temporal classification: labelling unsegmented sequence data with recurrent neural networks,'' in \textit{Proceedings of the 23rd International Conference on Machine Learning}, 2006, pp. 369-376.
    
    \item S. Watanabe et al., ``Hybrid CTC/Attention Architecture for End-to-End Speech Recognition,'' \textit{IEEE Journal of Selected Topics in Signal Processing}, vol. 11, no. 8, pp. 1240-1253, 2017.
    
    \item Z. Zhu, D. Li, and L. Wang, ``Visual Alignment Constraint for Continuous Sign Language Recognition,'' in \textit{Proceedings of the IEEE/CVF International Conference on Computer Vision (ICCV)}, 2021, pp. 11542-11551.
    
    \item A. Vaswani et al., ``Attention is all you need,'' in \textit{Advances in Neural Information Processing Systems}, vol. 30, 2017.
    
    \item K. He, X. Zhang, S. Ren, and J. Sun, ``Deep Residual Learning for Image Recognition,'' in \textit{Proceedings of the IEEE Conference on Computer Vision and Pattern Recognition (CVPR)}, 2016, pp. 770-778.
    
    \item D. P. Kingma and J. Ba, ``Adam: A Method for Stochastic Optimization,'' in \textit{International Conference on Learning Representations (ICLR)}, 2015.
\end{enumerate}

\vspace{1cm}

\section*{Acknowledgements}
\addcontentsline{toc}{section}{Acknowledgements}

I would like to express my sincere gratitude to my supervisor, Prajna P.N., for guidance and support throughout this project. I also thank IIT Patna for providing the computational resources necessary for training deep learning models.

Special thanks to the creators of the RWTH-PHOENIX-Weather 2014 dataset for making their data publicly available for research purposes.

\vspace{1cm}

\section*{Declaration}
\addcontentsline{toc}{section}{Declaration}

I hereby declare that the work presented in this report is my own original work carried out during Semester III of my MTech program at IIT Patna. I have not plagiarized or submitted this work elsewhere for any other degree or diploma.

\vspace{2cm}

\noindent
\begin{minipage}[t]{0.45\textwidth}
\textbf{Student Signature:} \hrulefill\\[0.5cm]
\textbf{Name:} Archie Narayan\\[0.3cm]
\textbf{Roll No:} 24A03RES119\\[0.3cm]
\textbf{Date:} December 2025
\end{minipage}
\hfill
\begin{minipage}[t]{0.45\textwidth}
\textbf{Supervisor Signature:} \hrulefill\\[0.5cm]
\textbf{Name:} \hrulefill\\[0.3cm]
\textbf{Designation:} \hrulefill\\[0.3cm]
\textbf{Date:} \hrulefill
\end{minipage}

\end{document}

